\documentclass[11pt]{amsart}
\usepackage[T1]{fontenc}
\usepackage[utf8]{inputenc}
\usepackage[french]{babel}
\usepackage{amssymb}
\usepackage[margin=1.05in]{geometry}
\usepackage[hidelinks]{hyperref}
\newcommand{\stamptime}{\the\numexpr\time/60\relax:\ifnum\numexpr\time-60*(\time/60)\relax<10 0\fi\the\numexpr\time-60*(\time/60)\relax}
\AddToHook{shipout/foreground}{\put(0.5\paperwidth,-\paperheight+0.55cm){\makebox[0pt]{\tiny\ttfamily unit/doc/\jobname.tex \ --- \ compiled \today, \stamptime}}}
\newtheorem{theorem}{Théorème}
\newtheorem{proposition}[theorem]{Proposition}
\newtheorem{lemma}[theorem]{Lemme}
\newtheorem{corollary}[theorem]{Corollaire}
\theoremstyle{definition}
\newtheorem{definition}[theorem]{Définition}
\theoremstyle{remark}
\newtheorem{remark}[theorem]{Remarque}
\DeclareMathOperator{\Syz}{Syz}
\DeclareMathOperator{\lc}{lc}
\DeclareMathOperator{\Ann}{Ann}
\newcommand{\sm}{\mathrm{span}}
\title{Artin--Rees constructif : preuve informelle et plan pour Unit}
\date{27 octobre 2026}
\begin{document}
\maketitle

\noindent\textbf{Statut.}
Cette note répond à \texttt{artin-rees-handoff.txt}. Elle fixe les
définitions et donne les preuves sur papier : noethérianité et
saturation (\S\S1--3), théorème de Hilbert pour les anneaux cohérents
fortement discrets noethériens (\S4), Artin--Rees (\S5), conséquences
(\S6), application à la platitude (\S7) et à ACL (\S8). La section~9
sépare ce qui est démontré sur papier, vérifié par Unit, et ouvert.
Toutes les preuves suivent Ruitenburg~[R], dans la variante \emph{avec
idéaux détachables}, qui est celle que [R, p.~48] donne sans choix
dépendant.

\section{Cadre}

Les anneaux sont commutatifs, unitaires, et des ensembles.
Pour $u_1,\dots,u_k\in R^r$ on note $\sm(u)$ le sous-module engendré.

\begin{definition}\label{def:cfn}
Un anneau $R$ est \emph{CFN} s'il vérifie :
\begin{itemize}
\item[(D)] \emph{forte discrétion} : pour $u_1,\dots,u_k,v\in R$, la
  proposition $v\in\sm(u)$ est décidable ;
\item[(C)] \emph{cohérence} : pour $u_1,\dots,u_k\in R$, il existe
  $\sigma^1,\dots,\sigma^p\in\Syz(u)=\{\lambda\in R^k\mid\sum\lambda_iu_i=0\}$
  tels que tout élément de $\Syz(u)$ soit combinaison des $\sigma^j$ ;
\item[(N)] \emph{noethérianité} au sens de Coquand--Persson : « Good bars
  $[\,]$ » (§\ref{sec:noeth}), pour une présentation de $R$ par des
  codes, surjective sous troncation (§\ref{sec:hilbert}).
\end{itemize}
\end{definition}

(D) donne l'égalité décidable ($v\in\sm(\,)$ signifie $v=0$).
Dans (D), seule la décision est une donnée ; le témoin $v=\sum\lambda_iu_i$
est donné sous troncation, comme les $\sigma^j$ de (C). Les trois
conditions sont donc des propositions, et CFN aussi. Les constructions
ci-dessous ont toutes une proposition pour but, et la récurrence se fait
dans ce but ; aucune suite infinie de données n'est choisie.

\begin{remark}
(D) est une hypothèse explicite, nommée et déchargée dans chaque
instance. Ruitenburg signale [R, p.~48] que, sans détachabilité, les
preuves connues de « une présentation finie sur un anneau cohérent
noethérien est cohérente noethérienne » et du théorème de Hilbert pour
les anneaux cohérents noethériens utilisent le choix dépendant. Les deux
applications ne demandent que des anneaux de présentation finie sur
$\mathbb Z$ ou sur un corps discret, qui vérifient (D).
\end{remark}

\section{Noethérianité et listes libres}\label{sec:noeth}

Soit $M$ un $R$-module. Une liste $s=[x_n,\dots,x_1]$ (la tête est
l'élément le plus récent) est \emph{bonne} (Good) si un $x_i$ est dans
$\sm(x_1,\dots,x_{i-1})$. $M$ est noethérien si Good bars $[\,]$ :
c'est la définition de \texttt{sag/noetherian} (\texttt{Span.Over.Noeth},
avec les éléments eux-mêmes pour codes, ou des codes surjectifs sous
troncation). La liste est \emph{libre} si aucun $x_i$ n'est dans le
sous-module engendré par les précédents.

\begin{lemma}[principe des listes libres]\label{lem:free}
Si $M$ est noethérien, la relation « $x::s$ prolonge $s$, avec $x\notin\sm(s)$ »
est bien fondée sur les listes libres : pour un prédicat propositionnel
$\Phi$ sur les listes, si $\Phi(x::s)$ pour tout $x\notin\sm(s)$ entraîne
$\Phi(s)$, alors $\Phi$ vaut sur toute liste libre.
\end{lemma}

\begin{proof}
Un arbre valide en $[\,]$ est valide, après descente le long de $s$, en
toute liste $s$. Récurrence sur l'arbre avec le prédicat
« $s$ libre $\Rightarrow\Phi(s)$ » : à une feuille, $s$ est bonne et libre,
contradiction ; à un n\oe ud, l'hypothèse de récurrence donne
$\Phi(x::s)$ pour tout $x$, en particulier pour $x\notin\sm(s)$, car
alors $x::s$ est libre.
\end{proof}

Avec (D), on reste dans les listes libres en n'ajoutant que des
non-membres, ce que (D) permet de décider. C'est la seule manière
dont la noethérianité est utilisée. La remarque du handoff
(\S3) est ainsi respectée : on n'utilise jamais une dépendance trouvée
au hasard, mais la bonne fondation des prolongements libres. On
n'affirme pas l'équivalence avec les variantes de Kuroki [K, 2.2--2.3].

\begin{proposition}[modules libres]\label{prop:free}
Si $R$ est CFN, alors pour tout $r$, le module $R^r$ vérifie (D), (C)
(pour les familles finies de $R^r$) et le principe du lemme~\ref{lem:free}.
\end{proposition}

\begin{proof}
Récurrence sur $r$, avec $R^r=R^{r-1}\oplus R$ et $\pi$ la dernière
coordonnée. Soient $s_1,\dots,s_k\in R^r$ et $N=\sm(s)$.

\emph{(C).} Par (C) dans $R$, des générateurs $\tau^1,\dots,\tau^p$ de
$\Syz(\pi s)$. Les vecteurs $w_j=\sum_i\tau^j_is_i$ sont dans $R^{r-1}$.
Par récurrence, des générateurs $\rho^1,\dots,\rho^q$ de $\Syz(w)$.
Alors $\Syz(s)$ est engendré par les $\sum_j\rho^\ell_j\tau^j$ : si
$\sum\lambda_is_i=0$, alors $\lambda=\sum c_j\tau^j$ et $\sum c_jw_j=0$,
donc $c=\sum d_\ell\rho^\ell$.

\emph{Intersection avec $R^{r-1}$.} $N_0=N\cap R^{r-1}$ est engendré par
les $w_j$ : si $\sum\lambda_is_i\in R^{r-1}$, alors $\lambda\in\Syz(\pi s)$.

\emph{(D).} Pour $v\in R^r$ : décider $\pi v\in\sm(\pi s)$. Si non,
$v\notin N$. Si oui, $\pi v=\sum\lambda_i\pi s_i$, et $v\in N$ si et
seulement si $v-\sum\lambda_is_i\in N_0$, ce qui se décide par
récurrence.

\emph{Bonne fondation.} Ordre lexicographique : on maintient une liste
libre $p$ de $R$ engendrant $\pi(N)$, et une liste libre $q$ de $R^{r-1}$
engendrant $N_0$. Soit $x\notin N$. Si $\pi x\notin\pi(N)$ (décidable),
$p$ est prolongée librement par $\pi x$ et $q$ est reconstruite à partir
de $[\,]$ pour le nouveau $N_0$ (en n'ajoutant que des non-membres, ce
qui termine par récurrence sur $r$). Sinon $x-y\in R^{r-1}$ pour un
$y\in N$, $p$ ne change pas et $q$ est prolongée librement par $x-y$
(qui n'est pas dans $N_0$ car $x\notin N$). On conclut par récurrence
sur $p$ (lemme~\ref{lem:free} dans $R$), puis sur $q$.
\end{proof}

\begin{corollary}\label{cor:inter}
Pour des familles finies de $R^r$ : intersection de deux sous-modules
de type fini, noyau d'une application $R^k\to R^r$, image réciproque
d'un sous-module de type fini, sont de type fini, avec générateurs
sous troncation.
\end{corollary}

\begin{proof}
$\sm(u)\cap\sm(v)$ est l'image de $\Syz(u_1,\dots,u_k,-v_1,\dots,-v_\ell)$
par la projection sur les $k$ premières coordonnées. Le reste est
analogue.
\end{proof}

\section{Saturation (lemme 2 de Ruitenburg)}\label{sec:sat}

Soit $M=R^r$ et $M[X]_n$ les polynômes de degré formel $\le n$ ; comme
$R$-module, $M[X]_n=R^{r(n+1)}$. Soient $f_1,\dots,f_s\in M[X]_n$ et
$N=R[X]f_1+\dots+R[X]f_s$.

\begin{lemma}\label{lem:step}
Pour un sous-module de type fini $P=\sm(p_1,\dots,p_t)$ de $M[X]_n$,
le sous-module $XP\cap M[X]_n$ est de type fini, explicitement.
\end{lemma}

\begin{proof}
$q=\sum\lambda_jp_j$ vérifie $Xq\in M[X]_n$ si et seulement si le
coefficient de degré $n$ de $q$ est nul, c'est-à-dire
$\lambda\in\Syz(c_n(p_1),\dots,c_n(p_t))$ dans $R^r$. Avec des
générateurs $\sigma^k$ (proposition~\ref{prop:free}), $XP\cap M[X]_n$
est engendré par les $X\sum_j\sigma^k_jp_j$.
\end{proof}

\begin{proposition}[saturation]\label{prop:sat}
Si $R$ est CFN, il existe (sous troncation) $p_1,\dots,p_t\in N\cap M[X]_n$,
avec leurs témoins dans $N$, tels que $P=\sm(p)$ contienne les $f_i$ et
vérifie $XP\cap M[X]_n\subseteq P$.
\end{proposition}

\begin{proof}
On part de $[\,]$ et on ajoute les $f_i$ en sautant les membres. Étape :
calculer les générateurs $g_1,\dots,g_u$ de $XP\cap M[X]_n$
(lemme~\ref{lem:step}) ; par (D) dans $M[X]_n$, ou bien tous sont dans
$P$ et on s'arrête, ou bien on ajoute le premier non-membre. La liste
reste libre, donc l'itération termine par la
proposition~\ref{prop:free} (principe des listes libres dans
$R^{r(n+1)}$). Chaque $g$ ajouté est dans $N$ : $g=Xq$ avec $q\in P\subseteq N$.
\end{proof}

\begin{theorem}[Ruitenburg, lemme 2]\label{thm:sat}
Avec $P$ comme ci-dessus, pour tout $m\ge n$,
\[N\cap M[X]_m=\textstyle\sum_{i=0}^{m-n}X^iP.\]
\end{theorem}

\begin{proof}
$\supseteq$ : $P\subseteq N$. $\subseteq$ : soit $f=\sum g_if_i\in M[X]_m$
avec $\deg g_i\le d$ ; récurrence sur $d$, pour tous $m\ge n$ et $f$
simultanément. Si $d=0$, $f\in\sm(f_i)\subseteq P$. Sinon
$g_i=g_i(0)+Xh_i$, $f^*=\sum h_if_i\in N$ et $f=Xf^*+\sum g_i(0)f_i$.
Comme $\sum g_i(0)f_i\in M[X]_n\subseteq M[X]_m$, on a $Xf^*\in M[X]_m$,
donc $f^*\in M[X]_{m-1}$ (et $f^*=0$ si $m=0$). Si $m=n$ : $f^*\in N\cap M[X]_n$,
donc $f^*\in P$ par récurrence, puis $Xf^*\in XP\cap M[X]_n\subseteq P$
et $f\in P$. Si $m>n$ : par récurrence $f^*\in\sum_{i\le m-1-n}X^iP$,
donc $f\in\sum_{i\le m-n}X^iP$.
\end{proof}

\begin{corollary}\label{cor:degree}
Pour tout $m$, $N\cap M[X]_m$ est un $R$-module de type fini,
engendré par les $X^ip_j$ ($i\le m-n$) si $m\ge n$, et par
$P\cap M[X]_m$ (corollaire~\ref{cor:inter}) si $m<n$. En particulier
$N\cap M$ est de type fini, et l'appartenance à $N$ est décidable,
avec témoin : $g\in M[X]_m$ est dans $N$ si et seulement s'il est
dans $\sum_{i\le m-n}X^iP$, ce que décide (D) dans $R^{r(m+1)}$.
\end{corollary}

\section{Le théorème de Hilbert pour les anneaux CFN}\label{sec:hilbert}

\begin{theorem}\label{thm:hilbert}
Si $R$ est CFN, alors $R[X]$ est CFN.
\end{theorem}

On travaille dans $S=R[X]$ sans choisir de représentants : la suite des
coefficients $\kappa:S\to R^{\mathbb N}$ est additive, injective, et
vérifie $\kappa(ag)=a\kappa(g)$ pour $a\in R$ et $\kappa(Xg)=X\kappa(g)$
(décalage). Tout $g$ s'écrit $g=g(0)+X\,\mathrm{tail}(g)$.

\emph{Pont.} Pour $u_1,\dots,u_k\in S$, soit $N$ le $R[X]$-module
engendré par les suites $\kappa(u_i)$ (forme récursive de
\S\ref{sec:sat}). Alors $y\in\sum_iSu_i$ si et seulement si
$\kappa(y)\in N$. En effet, si $c\in S$, $\kappa(cu_i)\in N$ par
récurrence sur un représentant de $c$ ($c=c(0)+X\,\mathrm{tail}(c)$) ;
réciproquement, si $\kappa(y)=\sum a_i\kappa(u_i)+X g'$ avec $g'$ de
témoin de degré $d$, par récurrence $g'=\kappa(\sum c'_iu_i)$ et
$y=\sum(a_i+Xc'_i)u_i$ par injectivité de $\kappa$.

\emph{(D).} Sous troncation, une borne commune des degrés des $u_i$ et
de $y$ ; le corollaire~\ref{cor:degree} (pour $M=R$) décide
$\kappa(y)\in N$.

\emph{(N).} $S$ est un quotient d'un type de listes ; une barre sur ses
éléments ne peut pas se construire sans choix. On porte donc la
noethérianité par des \emph{codes} : une présentation
(\texttt{Span.Pres}) noethérienne au sens de Coquand--Persson, dont
l'application vers l'anneau est surjective sous troncation. Le théorème
de Hilbert de Coquand--Persson (\texttt{noetherian/Hilbert.seg}) donne
une telle présentation pour $S$ (codes polynomiaux). Le principe des
listes libres (lemme~\ref{lem:free}), seul usage de (N), se transporte
des codes aux éléments : une liste finie d'éléments a, sous troncation,
une liste de codes (choix fini), une liste libre d'éléments est libre en
codes, et le prédicat est propositionnel.

\emph{(C).} Soient $u_1,\dots,u_k\in S$ de degrés $\le n$. La
proposition~\ref{prop:sat} (pour $M=R$) donne $P=(p_1,\dots,p_t)$ dans
$R[X]_n$ avec $u_i=\sum_jC_{ij}p_j$ ($C_{ij}\in R$), $p_j=\sum_iB_{ji}u_i$
($B_{ji}\in S$, choix fini sous troncation) et $XP\cap R[X]_n\subseteq P$.
Soient $\rho^1,\dots,\rho^a$ des générateurs de $\Syz_R(p)$ ($p_j$ vu
dans $R^{n+1}$, proposition~\ref{prop:free}) et $\sigma^1,\dots,\sigma^b$
des générateurs de $\Syz_R(c_n(p_1),\dots,c_n(p_t))$. Le polynôme
$w_k=\sum_j\sigma^k_jp_j$ est de degré $<n$, donc
$Xw_k\in XP\cap R[X]_n\subseteq P$ : $Xw_k=\sum_jc^k_jp_j$. Posons
$\tau^k=X\sigma^k-c^k\in\Syz_S(p)$.

\begin{lemma}[syzygies de Schreyer]\label{lem:crit}
$\Syz_S(p)$ est engendré par les $\rho^i$ et les $\tau^k$.
\end{lemma}

\begin{proof}
Soit $\gamma\in\Syz_S(p)$ avec $\deg\gamma_j\le e$. Récurrence sur $e$.
Le coefficient de degré $n+e$ de $\sum_j\gamma_jp_j=0$ est
$\sum_j\gamma_{j,e}\,c_n(p_j)$ (car $\deg p_j\le n$), donc
$\gamma_e=(\gamma_{j,e})_j\in\Syz_R(c_n(p))$. Si $e=0$, $\gamma$ est
constant et $\gamma\in\Syz_R(p)$ est combinaison des $\rho^i$. Si $e>0$,
$\gamma_e=\sum_k\mu_k\sigma^k$, et
$\gamma'=\gamma-\sum_k\mu_kX^{e-1}\tau^k$ est une syzygie de degré
$\le e-1$ : son coefficient de degré $e$ est
$\gamma_e-\sum_k\mu_k\sigma^k=0$, car les $c^k$ sont constants. On
conclut par récurrence.
\end{proof}

Si $\lambda\in\Syz_S(u)$, alors $C^{\mathsf T}\lambda\in\Syz_S(p)$, donc
$C^{\mathsf T}\lambda=\sum_\ell a_\ell\theta^\ell$ pour les générateurs
$\theta^\ell$ du lemme, et
\[\lambda=\textstyle\sum_i\lambda_i\,(e_i-B^{\mathsf T}C^{\mathsf T}e_i)
+\sum_\ell a_\ell\,B^{\mathsf T}\theta^\ell .\]
Les $e_i-B^{\mathsf T}C^{\mathsf T}e_i$ et les $B^{\mathsf T}\theta^\ell$
sont des syzygies de $u$ :
$\sum_l(B^{\mathsf T}\theta)_lu_l=\sum_j\theta_jp_j=0$ et
$\sum_l(B^{\mathsf T}C^{\mathsf T}e_i)_lu_l=\sum_jC_{ij}p_j=u_i$. Ils
engendrent donc $\Syz_S(u)$.\qed


\begin{proposition}[quotients et localisations]\label{prop:closure}
Si $R$ est CFN, $R/(c_1,\dots,c_m)$ et $R[1/c]$ sont CFN.
\end{proposition}

\begin{proof}
\emph{Quotient.} (D) : $\bar v\in\sm(\bar u)$ ssi $v\in\sm(u,c)$. (C) :
projeter sur les $k$ premières coordonnées les générateurs de
$\Syz(u_1,\dots,u_k,c_1,\dots,c_m)$. (N) : \texttt{Transfer.Over.noeth}.

\emph{Localisation.} Soit $J=\sm(u)$ (dénominateurs chassés). Les idéaux
$J:c^k$ sont de type fini (corollaire~\ref{cor:inter}) et croissants, et
$J:c^k=J:c^{k+1}$ entraîne $J:c^{k+1}=J:c^{k+2}$ : si $c^{k+2}x\in J$,
alors $cx\in J:c^{k+1}=J:c^k$. Par (D) et le principe des listes libres,
on trouve $k$ avec $J:c^k=J:c^{k+1}$ (on ajoute un générateur de
$J:c^{k+1}$ qui n'est pas dans $J:c^k$, s'il y en a). Alors $v\in JR[1/c]$
ssi $v\in J:c^k$, décidable. (C) : $\Syz_{R[1/c]}(u)$ est engendré par
$\Syz_R(u)$, car $\sum s_iu_i=0$ dans $R[1/c]$ donne
$c^b\sum s_iu_i=0$ dans $R$. (N) : \texttt{Hilbert.Over} avec $S=R[1/c]$,
$X=1/c$.
\end{proof}

\begin{proposition}[instances de base]\label{prop:base}
$\mathbb Z$ et les corps discrets sont CFN.
\end{proposition}

\begin{proof}
Corps : (D) et (C) par élimination ; (N) : \texttt{Field.Over.noeth}.
$\mathbb Z$ : (D) par division euclidienne et Bézout ; (C) : par
Euclide étendu, une matrice unimodulaire $U$ envoie $(u_1,\dots,u_k)$ sur
$(g,0,\dots,0)$ ; $\Syz(u)$ est engendré par les images par $U$ des
$e_2,\dots,e_k$ (et de $e_1$ si $g=0$) ; (N) : \texttt{Integers.Over.noeth}.
\end{proof}

\section{Artin--Rees}\label{sec:ar}

\begin{proposition}[Artin--Rees, cas principal]\label{prop:arp}
Soient $R$ CFN, $J=(a_1,\dots,a_k)$ et $f\in R$. Il existe (sous
troncation) $c$ tel que, pour tout $n$ et tout $x\in R$,
\[x\in f^{n+c}R\cap J\ \Longrightarrow\ x=f^ny\ \text{avec}\ y\in f^cR\cap J,\]
c'est-à-dire $f^{n+c}R\cap J=f^n(f^cR\cap J)$ ; l'inclusion
$f^n(f^cR\cap J)\subseteq f^{n+c}R\cap J$ est immédiate.
\end{proposition}

\begin{proof}
Soit $J:f^t=\{m\mid f^tm\in J\}$. Par le lemme de stabilisation de la
preuve de la proposition~\ref{prop:closure} (principe des listes libres
sur les générateurs de $J:f^{t+1}$, dernières coordonnées des syzygies
de $(a_1,\dots,a_k,f^{t+1})$), il existe $K$ avec
$J:f^{K+1}\subseteq J:f^K$, d'où $J:f^t\subseteq J:f^K$ pour tout $t$ :
si $f^{j+K+1}m\in J$, alors $fm\in J:f^{j+K}\subseteq J:f^K$ par
récurrence sur $j$, donc $m\in J:f^{K+1}\subseteq J:f^K$. On prend $c=K$.
Si $x=f^{n+K}m\in J$, alors $m\in J:f^{n+K}\subseteq J:f^K$, et
$y=f^Km$ convient : $y\in f^KR\cap J$ et $x=f^ny$.
\end{proof}

\begin{theorem}[Artin--Rees]\label{thm:ar}
Soient $R$ CFN, $I=(b_1,\dots,b_m)$, $M$ de présentation finie et
$j:N\to M$ injective, $N$ de présentation finie. Il existe (sous
troncation) $c$ tel que pour tout $m'\ge0$
\[I^{c+m'}M\cap j(N)=I^{m'}(I^cM\cap j(N)).\]
\end{theorem}

L'inclusion $\supseteq$ est immédiate. La borne $c$ est obtenue avant
toute quantification sur $m'$ et sur les éléments.

\emph{Réduction au cas libre.} $M=R^r/K$, $K$ de type fini, et soit
$N'\subseteq R^r$ l'image réciproque de $j(N)$, de type fini (générateurs
relevés plus ceux de $K$). Comme $K\subseteq N'$,
$(I^nR^r+K)\cap N'=(I^nR^r\cap N')+K$. Donc l'énoncé pour
$(R^r,N')$ donne l'énoncé pour $(M,j(N))$.

On traite $I=(a,b)$, qui suffit aux applications ; pour
$I=(b_1,\dots,b_m)$ la preuve est la même avec $m$ variables au lieu de
$U,V$. Soit $N=\sm(n_1,\dots,n_s)\subseteq M=R^r$. Les puissances $I^n$
sont celles de \texttt{IdealPowers} : $x\in I^{n+1}M$ si et seulement si
$x=ax_1+bx_2$ avec $x_1,x_2\in I^nM$ ; de même
$x\in I^{m+1}P$, pour un sous-module $P$, si $x=ax_1+bx_2$ avec
$x_1,x_2\in I^mP$, et $I^0P=P$.

\subsection{Pièces homogènes}
Un polynôme homogène de degré $n$ en $U,V$ à coefficients dans $M$,
$\sum_{j\le n}g_jU^{n-j}V^j$, est représenté par la suite
$g=(g_0,g_1,\dots)$ de $M$, nulle à partir de $n+1$ (on pose $U=1$). La
multiplication par $U$ laisse la suite inchangée (le degré augmente de
$1$), la multiplication par $V$ la décale ($Xg$, \texttt{SeqPoly.Xs}).
On pose
\[\theta_0(g)=g_0,\qquad \theta_{n+1}(g)=a\,\theta_n(g)+b^{n+1}g_{n+1},\]
de sorte que $\theta_n(g)=\sum_{j\le n}a^{n-j}b^jg_j$, et
\[Q_n=\{g\mid g\text{ nulle à partir de }n+1,\ \theta_n(g)\in N\}.\]

\begin{lemma}\label{lem:theta}
(i) $\theta_n$ est $R$-linéaire. (ii) Si $g_{n+1}=0$,
$\theta_{n+1}(g)=a\,\theta_n(g)$. (iii) $\theta_{n+1}(Xq)=b\,\theta_n(q)$.
(iv) $x\in I^nM$ si et seulement si $x=\theta_n(g)$ pour une suite $g$
nulle à partir de $n+1$. (v) $\theta_{n+1}(g)=a^{n+1}g_0+b\,\theta_n(g')$,
où $g'_j=g_{j+1}$. En particulier $Q_n$ est un sous-module,
$Q_n\subseteq Q_{n+1}$ et $XQ_n\subseteq Q_{n+1}$.
\end{lemma}

\begin{proof}
(iii) par récurrence sur $n$ :
$\theta_1(Xq)=a\cdot0+bq_0$ et
$\theta_{n+2}(Xq)=a\,b\,\theta_n(q)+b^{n+2}q_{n+1}=b\,\theta_{n+1}(q)$.
(iv) Si $x=ax_1+bx_2$ avec $x_i=\theta_n(g^i)$, alors
$x=\theta_{n+1}(g^1+Xg^2)$ par (ii) et (iii) ; réciproquement
$\theta_{n+1}(g)=a\,\theta_n(g)+b\,(b^ng_{n+1})$ et $b^ny\in I^nM$.
(v) par récurrence sur $n$ :
$\theta_{n+2}(g)=a(a^{n+1}g_0+b\,\theta_n(g'))+b^{n+2}g_{n+2}
=a^{n+2}g_0+b\,\theta_{n+1}(g')$.
\end{proof}

Les décompositions $g=p+Xq$ avec $p,q\in Q_n$ forment un sous-module :
on les ajoute et on les multiplie par les constantes composante par
composante ; $g=p+X0$ si $g=p\in Q_n$, et $g=0+Xq$ si $g=Xq$, $q\in Q_n$.

\begin{definition}
$\mathrm{Split}(c)$ : pour tout $n\ge c$ et tout $g\in Q_{n+1}$, il existe
(sous troncation) $p,q\in Q_n$ avec $g=p+Xq$.
\end{definition}

\begin{lemma}\label{lem:split-ar}
Si $\mathrm{Split}(c)$, alors pour tout $m$ et tout
$x\in I^{m+c}M\cap N$, on a $x\in I^m(I^cM\cap N)$ (sous troncation).
\end{lemma}

\begin{proof}
Par le lemme~\ref{lem:theta} (iv), $x=\theta_{m+c}(g)$ avec
$g\in Q_{m+c}$. Récurrence sur $m$, pour tout $g\in Q_{m+c}$. Si $m=0$,
$\theta_c(g)\in I^cM\cap N$. Sinon $g=p+Xq$ avec
$p,q\in Q_{m-1+c}$, et
$\theta_{m+c}(g)=a\,\theta_{m-1+c}(p)+b\,\theta_{m-1+c}(q)$ (lemme~\ref{lem:theta}),
avec $\theta(p),\theta(q)\in I^{m-1}(I^cM\cap N)$ par récurrence.
\end{proof}

Il reste à trouver $c$ avec $\mathrm{Split}(c)$. Ce n'est pas une
stabilisation de chaîne (le handoff, \S3, le signale) : on l'obtient
d'une famille génératrice finie du module de Rees, calculée par le
corollaire~\ref{cor:degree}.

\subsection{L'anneau $R'=R[U][V]$ et les composantes homogènes}
$R'=R[U][V]$ (deux extensions polynomiales de la bibliothèque) est CFN
(théorème~\ref{thm:hilbert} deux fois). Pour $G\in R'^r$, soit
$\pi_n(G)$ la suite des coefficients de $U^{n-j}V^j$ dans $G$ ($j\le n$),
définie par récurrence à partir des coefficients en $V$ puis en $U$
(\texttt{PolyCoef.kap} deux fois) :
$\pi_0(G)_0=[U^0V^0]G$, $\pi_{n+1}(G)_0=[U^{n+1}V^0]G$,
$\pi_0(G)_{j+1}=0$, $\pi_{n+1}(G)_{j+1}=\pi_n(G')_j$ où $G'$ a pour
coefficients en $V$ ceux de $G$ décalés. On le fait une fois pour toutes
sur les suites doubles $D_{ij}$ (le coefficient de $U^iV^j$), puis
$\pi_n(G)=\pi_n(D(G))$, $D(G)_{ij}$ étant le vecteur des coefficients de
$U^iV^j$ des coordonnées de $G$.

\begin{lemma}\label{lem:pi}
(i) $\pi_n$ est additive et $\pi_n(cG)=c\,\pi_n(G)$ pour $c\in R$ ;
$\pi_n(G)$ est nulle à partir de $n+1$.
(ii) $\pi_0(UG)=0$, $\pi_{n+1}(UG)=\pi_n(G)$.
(iii) $\pi_0(VG)=0$, $\pi_{n+1}(VG)=X\pi_n(G)$.
(iv) Pour $m\in M$ (constant), $\pi_0(m)=m$ et $\pi_{n+1}(m)=0$.
\end{lemma}

\begin{proof}
Sur les suites doubles : (ii) par récurrence sur $n$, (iii) et (iv) par
définition. Puis $D(UG)=U\,D(G)$, $D(VG)=V\,D(G)$, $D(cG)=c\,D(G)$ :
\texttt{kap-cmul}, \texttt{kap-X} en $V$ puis en $U$.
\end{proof}

\begin{lemma}[récurrence sur $R'$]\label{lem:rind}
Soit $P$ un prédicat propositionnel sur un $R'$-module, vrai en $0$,
stable par addition, par multiplication par les constantes de $R$, par
$U$ et par $V$. Alors $P$ est stable par multiplication par tout
$h\in R'$.
\end{lemma}

\begin{proof}
Sous troncation, $h$ a une borne de degré en $V$ ; récurrence sur
celle-ci avec $h=h_0+V\,\mathrm{tail}(h)$, $h_0\in R[U]$ (\texttt{horner}),
puis, pour $h_0$, récurrence sur une borne en $U$ avec
$h_0=c+U\,\mathrm{tail}(h_0)$.
\end{proof}

\begin{lemma}[degré total]\label{lem:tot}
Tout $G\in R'^r$ a, sous troncation, un degré total $c$ : $\pi_n(G)=0$
pour $n\ge c$.
\end{lemma}

\begin{proof}
La propriété « avoir, sous troncation, un degré total » est vraie en $0$,
stable par somme (somme des bornes), par les constantes, par $U$ et par
$V$ (la borne augmente de $1$, lemme~\ref{lem:pi}), et vraie pour les
constantes $m\in M$ (borne $1$). Par le lemme~\ref{lem:rind} elle vaut
pour les $h\,e_k$, $h\in R'$, $e_k$ la base standard, et
$G=\sum_kG_ke_k$. Aucun choix fini n'est utilisé.
\end{proof}

On pose, pour une suite $g$ de $M$,
$\iota_0(g)=g_0$ et $\iota_{n+1}(g)=U^{n+1}g_0+V\iota_n(g')$ dans
$R'^r$ ($g'_j=g_{j+1}$), de sorte que $\iota_n(g)=\sum_{j\le n}g_jU^{n-j}V^j$
et $\pi_n(\iota_n(g))=g$ si $g$ est nulle à partir de $n+1$ : par
récurrence sur $n$, $\pi_{n+1}(U^{n+1}g_0)=\pi_0(g_0)=(g_0,0,\dots)$ et
$\pi_{n+1}(V\iota_n(g'))=X\pi_n(\iota_n(g'))=Xg'$ (lemme~\ref{lem:pi}).

\subsection{Le module de Rees par le corollaire~\ref{cor:degree}}
Dans $R'^r[Y]$ (suites de $R'^r$, comme au \S\ref{sec:sat}), soit $Q^+$ le
sous-$R'[Y]$-module engendré par les $n_k$ (constants, de degré $0$) et
les $(U-aY)e_k$, $(V-bY)e_k$ ($k\le r$, de degré $1$). Par le
corollaire~\ref{cor:degree} pour l'anneau CFN $R'$, $Q^+\cap R'^r$ (degré
$0$ en $Y$) est engendré, sous troncation, par des
$G_1,\dots,G_p\in R'^r$. Soit $c$ un majorant de leurs degrés totaux
(lemme~\ref{lem:tot}).

\begin{lemma}[bon]\label{lem:good}
Soit $\varepsilon_n(H)=\theta_n(\pi_n(H))$ pour $H\in R'^r$, et pour
$G=(G_k)_k\in R'^r[Y]$, par récurrence sur $t$,
$E(G)_0=\varepsilon_0(G_0)$, $E(G)_{t+1}=\varepsilon_{t+1}(G_0)+E(G')_t$
($G'_k=G_{k+1}$), de sorte que
$E(G)_t=\sum_{k\le t}\theta_{t-k}(\pi_{t-k}(G_k))$ (le coefficient de
$Y^t$ après $U\mapsto aY$, $V\mapsto bY$) ; $G$ est \emph{bon} si
$E(G)_t\in N$ pour tout $t$. Les éléments de $Q^+$ sont bons. En
particulier, si $G\in Q^+\cap R'^r$, alors $\pi_n(G)\in Q_n$ pour tout $n$.
\end{lemma}

\begin{proof}
$\varepsilon_{n+1}(UH)=a\,\varepsilon_n(H)$,
$\varepsilon_{n+1}(VH)=b\,\varepsilon_n(H)$,
$\varepsilon_0(UH)=\varepsilon_0(VH)=0$ (lemmes~\ref{lem:theta}
(ii), (iii), \ref{lem:pi}) ; d'où, par récurrence sur $t$,
$E(UG)_{t+1}=a\,E(G)_t$, $E(VG)_{t+1}=b\,E(G)_t$, $E(UG)_0=E(VG)_0=0$,
et $E(YG)_{t+1}=E(G)_t$, $E$ est
linéaire : être bon est stable par $U$, $V$, $Y$, les constantes, la
somme, donc (lemme~\ref{lem:rind}) par $R'[Y]$. Les générateurs sont
bons : $E(n_k)=(n_k,0,\dots)$, et $E((U-aY)e_k)_1=ae_k-ae_k=0$, de même
pour $V$, les autres termes étant nuls. Pour $G$ de degré $0$ en $Y$,
$E(G)_n=\theta_n(\pi_n(G))$.
\end{proof}

\begin{lemma}[télescopage]\label{lem:tele}
Pour $g\in Q_n$, $\iota_n(g)\in Q^+$.
\end{lemma}

\begin{proof}
Notons $x\sim y$ pour $x-y\in Q^+$. Pour $(h,c)=(U,a)$ ou $(V,b)$,
$\Phi(w)=h\,w-Y\,c\,w$ est $R'$-linéaire en $w\in R'^r$ et
$\Phi(e_k)=(h-cY)e_k$ est un générateur ; donc $h\,w\sim Y\,c\,w$ pour
tout $w$, en écrivant $w=\sum_kw_ke_k$. Par récurrence sur $k$,
$U^kx\sim Y^ka^kx$ pour $x\in M$ :
$U^{k+1}x\sim UY^ka^kx=Y^kU\,a^kx\sim Y^{k+1}a^{k+1}x$. Par récurrence sur
$n$, pour toute suite $g$, $\iota_n(g)\sim Y^n\theta_n(g)$ : par
définition de $\iota$ et le lemme~\ref{lem:theta} (v),
\[\iota_{n+1}(g)=U^{n+1}g_0+V\iota_n(g'),\qquad
Y^{n+1}\theta_{n+1}(g)=Y^{n+1}a^{n+1}g_0+Y^{n+1}b\,\theta_n(g'),\]
$U^{n+1}g_0\sim Y^{n+1}a^{n+1}g_0$ et
$V\iota_n(g')\sim VY^n\theta_n(g')=Y^nV\theta_n(g')\sim Y^{n+1}b\,\theta_n(g')$.
Si de plus $\theta_n(g)\in N$, $Y^n\theta_n(g)$ est une combinaison (à
coefficients constants) des $Y^nn_k$, donc $\iota_n(g)\in Q^+$.
\end{proof}

\begin{proposition}\label{prop:split}
$\mathrm{Split}(c)$.
\end{proposition}

\begin{proof}
Soit $S(G)$ : pour tout $n$, $\pi_n(G)\in Q_n$, et pour tout $n\ge c$,
$\pi_{n+1}(G)=p+Xq$ avec $p,q\in Q_n$ (sous troncation). $S$ est vrai en
$0$, stable par somme et par les constantes ; par $U$ :
$\pi_{n+1}(UG)=\pi_n(G)+X0$ ; par $V$ : $\pi_{n+1}(VG)=0+X\pi_n(G)$
(lemme~\ref{lem:pi}, et $Q_n\subseteq Q_{n+1}$, $XQ_n\subseteq Q_{n+1}$).
Les $G_i$ vérifient $S$ : $\pi_n(G_i)\in Q_n$ (lemme~\ref{lem:good}) et
$\pi_{n+1}(G_i)=0$ pour $n\ge c$. Par le lemme~\ref{lem:rind}, $S$ est
stable par $R'$, donc vaut sur $Q^+\cap R'^r$. Si $g\in Q_{n+1}$ avec $n\ge c$, $\iota_{n+1}(g)\in Q^+\cap R'^r$
(lemme~\ref{lem:tele}) et $\pi_{n+1}(\iota_{n+1}(g))=g$.
\end{proof}

Le théorème suit des lemmes~\ref{lem:split-ar} et de la
proposition~\ref{prop:split} ; l'inclusion $\supseteq$ est immédiate.
Seule la direction facile de la division par $U-aY$, $V-bY$ est utilisée
(lemme~\ref{lem:tele}) : le noyau de la substitution n'est pas calculé,
et aucune formule de produit de polynômes n'est nécessaire, seulement
les multiplications par $U$, $V$, $Y$ et les constantes.\qed

\subsection{Formalisation}
Vérifié par Unit (\texttt{sag/artin-rees}, 27 octobre) pour $M=R^r$ et
$N=\sm(n_1,\dots,n_s)$ ; la réduction du cas de présentation finie au cas
libre n'est pas formalisée.

\begin{tabular}{ll}
\texttt{ReesVec} & $\theta_n$, $Q_n$, lemme~\ref{lem:theta}, Split,
  lemme~\ref{lem:split-ar} (modules sur $R$ seulement)\\
\texttt{Bigraded} & $R'=R[U][V]$, $\pi_n$, $\iota_n$,
  lemmes~\ref{lem:pi}, \ref{lem:rind} et \ref{lem:tot}\\
\texttt{ReesPlus} & $Q^+$ dans le modèle de \texttt{DegreeCor}, $E$,
  lemmes~\ref{lem:good} et \ref{lem:tele}\\
\texttt{ArtinRees} & proposition~\ref{prop:split} et le théorème
\end{tabular}

\section{Conséquences}\label{sec:cons}

\begin{corollary}[relèvement matriciel]\label{cor:lift}
Soient $C_0,C_1$ de présentation finie et $d:C_0\to C_1$. Il existe $c$
tel que : si $dz\in I^{c+m}C_1$, il existe $w\in I^mC_0$ avec $dw=dz$.
\end{corollary}

\begin{proof}
$\operatorname{im}d$ est de type fini et de présentation finie
(cohérence). Artin--Rees pour $\operatorname{im}d\subseteq C_1$ donne
$dz\in I^m(I^cC_1\cap\operatorname{im}d)$, donc
$dz=\sum\beta_\ell\,dz_\ell$ avec $\beta_\ell\in I^m$ ; on pose
$w=\sum\beta_\ell z_\ell$.
\end{proof}

\begin{corollary}[coefficients]\label{cor:coef}
Soient $I=(a,b)$, $N\subseteq R^r$ engendré par $u_1,\dots,u_p$ et $c$ une
borne d'Artin--Rees pour $N$. Si $x\in I^{c+m}R^r\cap N$, alors il existe
$\lambda_1,\dots,\lambda_p\in I^m$ avec $x=\sum_j\lambda_ju_j$.
\end{corollary}

\begin{proof}
Le théorème~\ref{thm:ar} donne $x\in I^m(I^cR^r\cap N)$, un arbre de
profondeur $m$ dont les feuilles sont dans $N$. Récurrence sur $m$ : une
feuille est une combinaison des $u_j$ (coefficients dans $I^0=R$) ; à un
nœud $x=ax_1+bx_2$, si $x_i=\sum_j\lambda_{i,j}u_j$ avec
$\lambda_{i,j}\in I^m$, alors $x=\sum_j(a\lambda_{1,j}+b\lambda_{2,j})u_j$
et $a\lambda_{1,j}+b\lambda_{2,j}\in I^{m+1}$ par définition de $I^{m+1}$.
La conclusion est une proposition : les décompositions des feuilles,
données sous troncation, s'éliminent.
\end{proof}

\begin{corollary}[Kill uniforme]\label{cor:kill}
Soit $E$ CFN, dont les éléments forment un ensemble, et $f\in E$. Il
existe $c$ tel que $\Ann(f)\cap f^{c+1}E=0$. En particulier
\texttt{StableAnnihilator.Stable} $f\ (c+1)$ :
$f^{c+2}v=0\Rightarrow f^{c+1}v=0$.
\end{corollary}

\begin{proof}
Par (C) pour la famille à un élément $(f)$, $\Ann(f)=(g_1,\dots,g_p)$, les
$g_j$ étant les syzygies de $(f)$ : $g_jf=0$, et $xf=0$ entraîne
$x\in(g)$. La proposition~\ref{prop:arp} avec $J=\Ann(f)$ et $n=1$ :
si $x\in f^{1+c}E\cap\Ann(f)$, alors $x=fy$ avec $y\in\Ann(f)$, donc
$x=0$. Si $f^{c+2}v=0$, alors $x=f^{c+1}v$ est dans
$\Ann(f)\cap f^{c+1}E$, donc nul.
\end{proof}

\section{Platitude sur un anneau arbitraire}\label{sec:flat}

On instancie \texttt{Noetheriannote.Noetheriannote}, puis
\texttt{EtaleFlat2.Arbitrary}.

\emph{Calcul (composante N).} $\mathrm{Noeth}(A)$ := $A$ est CFN, énoncé
sur les éléments de \texttt{A.R} (une proposition, donc stable par les
quotients de la bibliothèque). Les quatre champs :
\texttt{polynomial} est le théorème~\ref{thm:hilbert} transporté sur
\texttt{PolynomialRing.Extension} (représentants sous troncation) ;
\texttt{quotient} et \texttt{localization} sont la
proposition~\ref{prop:closure} sur \texttt{QuotientRing} et
\texttt{PowerLocalization} ; \texttt{stable} est le
corollaire~\ref{cor:kill}. La platitude finale dépend donc d'Artin--Rees.

\emph{Réduction (composante D).} Données : $A$, $a_1,\dots,a_n$, $F,G$,
$p_i,u_F,u_G$ et l'identité (D1) dans $A[X][Y]$. Sous troncation, on
choisit des représentants (listes de coefficients) ; soit
$c_1,\dots,c_N\in A$ la liste finie de tous les coefficients qui y
figurent. Soit $A_0'=\mathbb Z[x_1,\dots,x_N]$ (itération de
\texttt{PolynomialRing.Extension} sur l'anneau des entiers) et
$h':A_0'\to A$, $x_j\mapsto c_j$. Les données se relèvent à $A_0'$ en
remplaçant $c_j$ par $x_j$. Les deux membres de (D1), calculés dans
$A_0'[X][Y]$, ont un nombre fini de coefficients ; leurs différences
$\delta_1,\dots,\delta_K\in A_0'$ sont envoyées sur $0$ par $h'$. On pose
$A_0=A_0'/(\delta_1,\dots,\delta_K)$ (\texttt{QuotientRing} de
\texttt{RegularIdeal.Generated}), et $h:A_0\to A$ induit. Dans $A_0$,
(D1) est vraie coefficient par coefficient ; le jacobien commute au
changement de base (\texttt{BaseChange}). $A_0$ est CFN par les
propositions~\ref{prop:base}, \ref{prop:closure} et le
théorème~\ref{thm:hilbert}. $h$ n'est pas injectif en général, ce que
l'interface permet.

Le piège signalé par le handoff ne se présente pas : on n'utilise pas le
sous-anneau image $E'$ de \texttt{RegularCoded}, mais un quotient de
présentation finie, dont la cohérence est construite.

\section{Application ACL}\label{sec:acl}

Anneau $A=k[b][a]$, $k$ corps discret, $I=(a,b)$ ; $A$ est CFN
(proposition~\ref{prop:base}, théorème~\ref{thm:hilbert} deux fois).

\subsection{Présentation (HP)}
Soit $S=A[X,Y,Z]$, $J=(P,Q)$, $\deg P=d$, $\deg Q=e$, $D_0=\gcd(P,Q)$ de
degré $\delta$, $P=D_0P'$, $Q=D_0Q'$, $t=d+e-\delta$. La suite
\[0\to S(-t)\xrightarrow{(-Q',P')}S(-d)\oplus S(-e)\xrightarrow{(P,Q)}J\to0\]
est exacte : si $\alpha P+\beta Q=0$ alors $\alpha P'=-\beta Q'$, et
$P',Q'$ premiers entre eux (données S0/S1 de \texttt{sag/gcd}) donnent
$\alpha=-\gamma Q'$, $\beta=\gamma P'$. Le complexe de Čech des trois
cartes est exact en chaque degré de torsion (localisations), d'où une
suite exacte longue. Pour le recouvrement standard,
$H^1(S(n))=0$, $H^3=0$ (pas de terme de degré 3), et $H^2(S(-s-1))$ est
libre de base les monômes $X^iY^jZ^k$ à exposants $<0$ de degré $-s-1$,
de rang $\binom s2$ ; en particulier $H^2(S(-1))=0$. En degré $-1$ :
$H^1(S/J)\cong H^2(J)$ (car $H^1(S(-1))=H^2(S(-1))=0$), et
\[H^2(S(-t-1))\to H^2(S(-d-1))\oplus H^2(S(-e-1))\to H^2(J(-1))\to0\]
est exacte. D'où une matrice finie
$L:A^{\binom t2}\to A^{\binom d2+\binom e2}$ de conoyau $H=H^1(C)$. Ses
colonnes : multiplier par $(-Q',P')$ puis appliquer l'opérateur $\pi$
(garder les monômes à exposants tous négatifs).

On n'utilise pas cette suite exacte abstraite : la réalisation suivante
dans le modèle de Čech de \texttt{sag/acl-zmt} en donne les deux moitiés
utiles (les colonnes de $L$ sont des relations ; toute relation en vient)
avec des témoins de bord explicites, sans définir $H$.

\subsection{Le modèle et les générateurs de (F1)}
Notations de \texttt{CechF1}, \texttt{CechContraction},
\texttt{CechSections} : les éléments de Laurent sont les fonctions
coefficients $\mathbb Z^3\to A$ ; $K^0(k)$, $K^1(k)$ les $0$- et
$1$-cochaînes de degré $k$ vérifiant les conditions de support des cartes
et à support fini ; $\delta_0,\delta_1$ ; $\pi$ (garder les monômes à trois
exposants $<0$), $\sigma$, $\tau$ avec
\[q=\pi(q)+\delta_1\sigma(q),\qquad \delta_0\tau(b)=b\ \text{ si }\delta_1b=0\]
($b$ vérifiant les supports), $\pi=0$ en degré $-1$, et
$\pi(\delta_1u)=0$ pour $u$ vérifiant les supports
(\texttt{CechSections.pi-d1}). Pour une cochaîne $u$ et un polynôme $F$
(exposants $\ge0$), $Fu$ vérifie encore les supports.

$P=\mu_P+aF_P+bG_P$, $Q=\mu_Q+aF_Q+bG_Q$ homogènes de degrés $d,e$ ;
$N^1=PK^1(-1-d)+QK^1(-1-e)$ (\texttt{CechOps.NW}, avec témoins) ;
$g=\mu_QP-\mu_PQ$, de degré $d_g=d+e$. Les générateurs de (F1) sont
$m_t=\sigma(gX^t)$, $t$ parcourant les $r=\binom{d_g}2$ exposants à trois
composantes $<0$ de degré $-d_g-1$, et $G\rho=\sum_t\rho_tm_t$. Pour
$\rho\in A^r$ on pose $q_\rho=\sum_t\rho_tX^t$ ; par linéarité de $\sigma$,
\[G\rho=\sigma(g\,q_\rho).\]
Le pgcd (\texttt{sag/gcd}, \texttt{HomLambda.Data}, $A=k[b][a]$) donne
$D_0$, $P'$, $Q'$ homogènes de degrés $\delta$, $d-\delta$, $e-\delta$,
avec $Pf=D_0P'f$, $Qf=D_0Q'f$, d'où $PQ'=QP'$, et
\begin{itemize}
\item[(S0)] $D_0f=0$, $f$ à support fini $\Rightarrow f=0$ ;
\item[(S1)] $P'\alpha+Q'\beta=0$, $\alpha,\beta$ homogènes à support fini
  $\Rightarrow$ il existe $w$, homogène, avec $\alpha=-Q'w$, $\beta=P'w$,
\end{itemize}
en tout degré (aucune hypothèse d'inversibilité de $D_0$). On note
$t_0=d+e-\delta$ et $s$ parcourt les $r'=\binom{t_0}2$ exposants à trois
composantes $<0$ de degré $-t_0-1$ ; $h_\eta=\sum_s\eta_sX^s$ pour
$\eta\in A^{r'}$.

\subsection{La matrice des relations}
\begin{definition}
$\mathcal M:A^r\times A^{r'}\to A^{\binom d2}\times A^{\binom e2}$,
\[\mathcal M(\rho,\eta)=\bigl(\pi(\mu_Qq_\rho-Q'h_\eta),\ \pi(P'h_\eta-\mu_Pq_\rho)\bigr),\]
lus sur les exposants à composantes $<0$ de degrés $-d-1$ et $-e-1$
(\texttt{CechMonomials.pi-gens} : $\pi(f)=\sum_uf(u)X^u$ pour $f$
homogène). Les degrés : $\mu_Qq_\rho$ et $Q'h_\eta$ sont de degré $-d-1$,
$\mu_Pq_\rho$ et $P'h_\eta$ de degré $-e-1$. Les coefficients de
$\mathcal M$ sont des coefficients de $\mu_P,\mu_Q,P',Q'$.
\end{definition}

\begin{lemma}[primitives]\label{lem:prim}
Si $\mathcal M(\rho,\eta)=0$, soient $\alpha=\mu_Qq_\rho-Q'h_\eta$,
$\beta=P'h_\eta-\mu_Pq_\rho$ et
\[b(\rho,\eta)=\tau\bigl(G\rho-P\sigma(\alpha)-Q\sigma(\beta)\bigr)\in K^0(-1).\]
Alors $G\rho-\delta_0b(\rho,\eta)=P\sigma(\alpha)+Q\sigma(\beta)\in N^1$.
\end{lemma}

\begin{proof}
$P\alpha+Q\beta=(\mu_QP-\mu_PQ)q_\rho+(QP'-PQ')h_\eta=g\,q_\rho$. Comme
$\pi(\alpha)=\pi(\beta)=0$, $\delta_1\sigma(\alpha)=\alpha$ et
$\delta_1\sigma(\beta)=\beta$ ; comme $gq_\rho$ est de degré $-1$,
$\delta_1\sigma(gq_\rho)=gq_\rho$. Donc l'argument $c$ de $\tau$ vérifie
$\delta_1c=gq_\rho-P\alpha-Q\beta=0$ ; il vérifie les supports, d'où
$\delta_0\tau(c)=c$ (contraction).
\end{proof}

\begin{proposition}[toute relation vient de $\mathcal M$]\label{prop:rel}
Si $G\rho-\delta_0v=PU+QV$ avec $v\in K^0(-1)$, $U\in K^1(-1-d)$,
$V\in K^1(-1-e)$, il existe $\eta$ avec $\mathcal M(\rho,\eta)=0$.
\end{proposition}

\begin{proof}
On applique $\delta_1$ : $\delta_1\delta_0=0$ et
$\delta_1G\rho=\delta_1\sigma(gq_\rho)=gq_\rho$ donnent
$gq_\rho=P\delta_1U+Q\delta_1V$, soit
\[P\alpha'+Q\beta'=0,\qquad \alpha'=\delta_1U-\mu_Qq_\rho,\ \ \beta'=\delta_1V+\mu_Pq_\rho .\]
Donc $D_0(P'\alpha'+Q'\beta')=0$, et (S0) donne $P'\alpha'+Q'\beta'=0$
(supports finis) ; (S1) donne $w$, de degré $-t_0-1$, avec
$\alpha'=-Q'w$, $\beta'=P'w$. On applique $\pi$. D'une part
$\pi(\delta_1U)=\pi(\delta_1V)=0$ ($U,V$ vérifient les supports). D'autre
part $w=\pi(w)+\delta_1\sigma(w)$ et $Q'\delta_1\sigma(w)=\delta_1(Q'\sigma(w))$,
où $Q'\sigma(w)$ vérifie les supports, d'où $\pi(Q'w)=\pi(Q'\pi(w))$, et
de même pour $P'$. Il vient
$\pi(\mu_Qq_\rho)=\pi(Q'\pi(w))$ et $\pi(\mu_Pq_\rho)=\pi(P'\pi(w))$ :
$\eta_s=w(s)$ convient, puisque $\pi(w)=h_\eta$ (\texttt{pi-gens}).
\end{proof}

Le lemme~\ref{lem:prim} dit que les colonnes de la matrice $L$ de (HP)
sont des relations, la proposition~\ref{prop:rel} que (HP) est exacte au
milieu ; seules ces deux moitiés servent.

\begin{corollary}[le module $N$]\label{cor:N}
Soit $N=\{\rho\mid\exists v,\ G\rho-\delta_0v\in N^1\}$. Il existe, sous
troncation, $(\rho_j,\eta_j)_{j\le p}$ engendrant $\ker\mathcal M$ (cohérence
de $A^{\binom d2+\binom e2}$, proposition~\ref{prop:free}). Avec
$u_j=\rho_j$ et $b_j=b(\rho_j,\eta_j)$ :
$Gu_j-\delta_0b_j\in N^1$, et $N=\sm(u_1,\dots,u_p)$.
\end{corollary}

\begin{proof}
$\supseteq$ : lemme~\ref{lem:prim} et linéarité. $\subseteq$ : pour
$\rho\in N$, la proposition~\ref{prop:rel} donne $\eta$ avec
$(\rho,\eta)\in\ker\mathcal M$, combinaison des $(\rho_j,\eta_j)$, donc
$\rho$ combinaison des $u_j$.
\end{proof}

\subsection{La correction}
$A=k[b][a]$ est CFN de base $K_0=k$ (proposition~\ref{prop:base},
théorème~\ref{thm:hilbert} deux fois). Les données de \texttt{CechF1} :
pour tout $n\ge1$, $\rho_n\in A^r$ à coordonnées dans $I^{n-1}$,
$c_n\in I^nK^0(-1)$ et
\[\delta_0z^*-(G\rho_n+\delta_0c_n)\in N^1\qquad(\beta_n),\]
donc $\rho_n\in N$ (avec $v=z^*-c_n$). Soit $c$ une borne d'Artin--Rees
pour $N=\sm(u_j)\subseteq A^r$ (théorème~\ref{thm:ar}), $m\ge0$ et
$n=c+m+1$. Alors $\rho_n\in I^{c+m}A^r\cap N$, et le
corollaire~\ref{cor:coef} donne $\rho_n=\sum_j\lambda_ju_j$ avec
$\lambda_j\in I^m$. Soit
\[w=c_n+\sum_j\lambda_jb_j\in I^mK^0(-1),\qquad h=z^*-w .\]
Alors $\delta_0w\equiv\delta_0c_n+\sum_j\lambda_jGu_j=\delta_0c_n+G\rho_n
\equiv\delta_0z^*$ modulo $N^1$, donc $\delta_0h\in N^1$, et
$h-z^*=-w\in I^mK^0(-1)$. Pour $m=1$ c'est la donnée
\texttt{ReesLiftingCong.Lift} ; tous les choix (générateurs, borne,
coefficients) sont sous troncation et s'éliminent dans la proposition
\texttt{TL}.

\subsection{Le certificat}
La suite est l'assemblage de \texttt{ACLZMTCong} (sections de degré~$0$
modulo $N^0$, restriction à la carte $Z$, raccourci (F2) de
\texttt{CechSections}), où la présentation noethérienne codée de $A$ servait
deux fois : pour le relèvement (remplacé par la correction ci-dessus) et
pour une dépendance entre les coordonnées des puissances d'une section
fermée (\texttt{GradedDependence} en l'idéal $(1,0)$). Celle-ci vient de
CFN :

\begin{lemma}[dépendance]\label{lem:dep}
Soit $M$ un module vérifiant (D) et le principe du
lemme~\ref{lem:free}, et $f_0,f_1,\dots\in M$. Il existe, sous troncation,
$N$ avec $f_N\in\sm(f_0,\dots,f_{N-1})$.
\end{lemma}

\begin{proof}
Soit $\Phi(s)$ : « pour tout $n$ avec $\sm(s)\subseteq\sm(f_0,\dots,f_{n-1})$,
la conclusion vaut » (une proposition). Soit $s$ libre avec $\Phi(x::s)$
pour tout $x\notin\sm(s)$, et $n$ comme dans $\Phi(s)$. On décide
$f_n\in\sm(s)$ par (D). Si oui, $f_n\in\sm(f_0,\dots,f_{n-1})$ : $N=n$.
Sinon $\Phi(f_n::s)$, et $\sm(f_n::s)\subseteq\sm(f_0,\dots,f_n)$ : on
l'applique en $n+1$. Le principe donne $\Phi([\,])$, appliqué en $n=0$.
\end{proof}

Pour $M=A^s$, $A$ CFN (proposition~\ref{prop:free}). Soit $c$ une section
fermée ; avec $f_n=\kappa(c^n)$ (coordonnées (F2)), le
lemme~\ref{lem:dep} donne $\kappa(c^N)=\sum_{i<N}x_i\kappa(c^i)$, donc
$c^N-\sum_ix_ic^i\in N^0$ et, après restriction,
$\mathrm{res}(c)^N=\sum_ix_i\,\mathrm{res}(c)^i$ : $\mathrm{res}(c)$ est
entier sur $A$. Avec $h$ donné par la correction ($m=1$), $Zh$, $Xh$, $Yh$
sont fermées, $s=\mathrm{res}(Zh)=1+a\,\mathrm{res}(Zm_1)+b\,\mathrm{res}(Zm_2)$
($h-z^*=am_1+bm_2$, $\mathrm{res}(Zz^*)=1$), $\mathrm{res}(Xh)=sx$,
$\mathrm{res}(Yh)=sy$ : $s$, $sx$, $sy$ sont entiers sur $A$. Les données
(F2) sont celles de \texttt{CechZMTD0} en degré~$0$ (données
$D_0,P',Q'$, (S0), (S1) de \texttt{HomLambda.Data}, avec
$k_W=-(e'+d_0+1)$).

\section{Ce qui est démontré, vérifié, ouvert}

\begin{itemize}
\item Sur papier, ici : \S\S1--8 ; pour \S8, dans le modèle de Čech des
  fichiers (lemme~\ref{lem:prim}, proposition~\ref{prop:rel},
  corollaire~\ref{cor:N}, la correction).
\item Vérifié par Unit (\texttt{sag/artin-rees/ACLRelations}, 27 octobre) :
  la matrice $\mathcal M$, le lemme~\ref{lem:prim}, la
  proposition~\ref{prop:rel} et le corollaire~\ref{cor:N}, dans le modèle
  de Čech de \texttt{sag/acl-zmt} (fichiers liés sans modification), sous
  les données $D_0,P',Q'$, (S0), (S1) de \texttt{HomLambda.Data} en
  degré $-1$ (comme \texttt{CechZMTD0}) et la cohérence de $A$. Les
  vecteurs de $\ker\mathcal M$ sont dans $A^{r+r'}$, $\mathcal M$ est
  linéaire et $\ker\mathcal M$ est l'ensemble des syzygies de ses colonnes
  $\mathcal M(e_k)$ (proposition~\ref{prop:free} pour
  $A^{\binom d2+\binom e2}$) ; $N$ est un sous-module (la troncation du
  témoin $v$), d'où $\sm(u_j)\subseteq N$.
\item Vérifié par Unit (\texttt{sag/artin-rees/ACLCorrection},
  \texttt{ACLCorrectionGeneral}, 27 octobre) : la correction, pour $A$
  CFN et tout $m$ (sous troncation, $h$ avec $\delta_0h\in N^1$ et
  $h-z^*\in I^mK^0(-1)$) ; pour $m=1$ c'est la donnée
  \texttt{ReesLiftingCong.Lift} des données (F1), sans présentation
  noethérienne codée de $A$. Les paramètres de \texttt{ACLRelations}
  viennent de \texttt{HomLambda.Data} ($\deg\mu_P=d_0+1$,
  $k_W=-(e'+d_0+2)$, (S1) transporté aux degrés $k_P,k_Q$).
\item Vérifié par Unit (\texttt{SeqDependence}, \texttt{ACLCertificate},
  \texttt{ACLSections}, \texttt{ACLCertificateGeneral}, 27 octobre) : le
  lemme~\ref{lem:dep} et le certificat ACL pour deux équations homogènes
  sur un anneau à pgcd $A$ CFN (instance de \texttt{HomLambda.Data}),
  sans présentation noethérienne codée ; la conclusion est celle de
  \texttt{ACLZMTCong}. \texttt{ACLSections} recopie la partie de
  \texttt{CechZMTD0} qui n'utilise pas la présentation (ce module la prend
  en paramètre global). L'instance $A=k[b][a]$, $k$ corps discret
  (\texttt{ACLCertificateK} : CFN par la proposition~\ref{prop:base} et le
  théorème~\ref{thm:hilbert} deux fois), sans autre hypothèse.
\item Vérifié par Unit avant ce travail : noethérianité CP de $\mathbb Z$,
  des corps, de $R[X]$, des quotients et localisations (codes) ;
  \texttt{Sequence.Kill} ; les données (F1) ; les interfaces
  \texttt{Noetheriannote} et \texttt{EtaleFlat2.Arbitrary}.
\item Vérifié par Unit (\texttt{sag/artin-rees}) : définitions CFN,
  proposition~\ref{prop:free}, transport le long d'isomorphismes
  linéaires ($M[X]_n\cong R^{r(n+1)}$), lemme~\ref{lem:step},
  proposition~\ref{prop:sat} et théorème~\ref{thm:sat} (polynômes
  représentés par des suites nulles à partir d'un rang ; $N$ par la forme
  récursive $f=\sum g_i(0)f_i+Xf^*$) ; corollaire~\ref{cor:degree}
  (pour tout $m$, y compris $m<n$, par le noyau de la troncation) ;
  théorème~\ref{thm:hilbert} : le pont $\kappa$, (D), le
  lemme~\ref{lem:crit} et (C), et (N) par les codes polynomiaux de
  Coquand--Persson.
\item Vérifié par Unit (quatrième session) : CFN est petit (codes
  $\mathrm{El}\,K_0\,c$, $c$ dans le type de données
  $\mathrm{CT}=\mathrm{base}\mid\mathrm{list}\,c$, base $K_0$ en
  paramètre ; « $R$ est CFN » est dans $U_l$) ; la
  proposition~\ref{prop:closure} (quotients, pour toute application
  surjective sous troncation de noyau $(c_1,\dots,c_m)$, et localisations,
  avec \texttt{QuotientRing} et \texttt{PowerLocalization} comme
  instances ; pour la localisation, la preuve formalisée passe par des
  dénominateurs communs, la détachabilité de $J:c^\infty$ et
  \texttt{Hilbert.Over} avec $X=1/c$ --- le texte ci-dessus est à
  réécrire en ce sens) ; la proposition~\ref{prop:base} pour les corps
  discrets ($1\neq0$).
\item Vérifié par Unit : proposition~\ref{prop:arp} ;
  théorème~\ref{thm:ar} pour $I=(a,b)$, $M=R^r$, $N$ de type fini
  (\texttt{ReesVec}, \texttt{Bigraded}, \texttt{ReesPlus},
  \texttt{ArtinRees}) ; Kill pour $I=(f)$ (\texttt{Kill}) ;
  le corollaire~\ref{cor:coef} (\texttt{PowSpan}) ; pour une base $K_0$
  fixée, les anneaux CFN forment un \texttt{Noetheriannote.Calculus}
  (\texttt{CFNCalculus} : \texttt{HilbertCFN}, \texttt{QuotCFNInst},
  \texttt{LocCFNInst}, \texttt{Kill.stable}), sans la composante
  \texttt{reduction}.
\item Vérifié par Unit (\texttt{ZRing}, \texttt{ZArith}, \texttt{ZCFN},
  27 octobre) : l'anneau $\mathbb Z$, sa propriété universelle (le morphisme
  $\mathbb Z\to A$ existe et est unique), et la proposition~\ref{prop:base}
  pour $\mathbb Z$ : (D) par la divisibilité par le pgcd (Euclide avec
  coefficients de Bézout), (C) par les générateurs $e_i-q_ib$
  ($u_i=q_iG$, $\sum b_iu_i=G$, simplification par $G\neq0$), (N) par
  l'arbre d'Euclide de \texttt{Integers}, codes $\mathbb N\times\mathbb N$.
\item À vérifier par Unit, dans l'ordre : la réduction de théorème~\ref{thm:ar} au cas libre ;
  corollaire~\ref{cor:lift} ; \S7. L'état courant est tenu dans
  \texttt{sag/artin-rees/README.txt}.
\item Taille : résolu par l'indexation des codes ; pour \S7,
  $\mathrm{Noeth}(A)$ := CFN de base les codes entiers $\mathbb N\times\mathbb N$.
\end{itemize}

\begin{thebibliography}{9}
\bibitem{R} W. Ruitenburg, Constructions of finitely generated submodules
  of constructively noetherian modules, \emph{Compositio Math.} 62 (1987),
  47--52.
\bibitem{CP} T. Coquand, H. Persson, Gröbner bases in type theory,
  TYPES'98, LNCS 1657 (1999).
\bibitem{K} R. Kuroki, A quantitative Hilbert's basis theorem and the
  constructive Krull dimension.
\bibitem{MRR} R. Mines, F. Richman, W. Ruitenburg, \emph{A course in
  constructive algebra}, Springer (1988).
\end{thebibliography}
\end{document}
